\chapter{РЕЗУЛЬТАТЫ И ИХ ОБСУЖДЕНИЕ}

\section{ВЫБОР И ОБОСНОВАНИЕ ОПТИМАЛЬНЫХ УСЛОВИЙ АНАЛИЗА}

Для достижения экспрессности, селективности и воспроизводимости результатов разделения неорганических анионов на приборе «Капель-105М» проведена оптимизация ключевых рабочих параметров.

\subsection{Выбор длины волны детектирования}

Отсутствие собственных хромофорных групп у большинства целевых анионов ($\text{Cl}^-$, $\text{SO}_4^{2-}$, $\text{F}^-$, $\text{HPO}_4^{2-}$) делает невозможным их прямое фотометрическое определение в ультрафиолетовой области спектра~\cite{chenUltravioletSpectroscopicDetection2021}. Для решения этой проблемы применен метод косвенного УФ-детектирования~\cite{pobozyApplicationCapillaryElectrophoresis2021b}. Метод основан на введении в состав фонового электролита специального ко-иона (визуализатора), обладающего высоким светопоглощением в выбранном оптическом диапазоне~\cite{komarova2006prakticheskoe}. В качестве визуализатора выбран хромат-анион ($\text{CrO}_4^{2-}$), имеющий интенсивную полосу поглощения с максимумом около 254~нм~\cite{shelekhovaOpredelenieAnionovNeorganicheskikh2023}. Детектирование проводили с использованием встроенной дейтериевой лампы при длине волны 254~нм~\cite{pndf14123428218}. Прохождение зоны неабсорбирующего аналита через окно детектора сопровождается локальным снижением концентрации поглощающего хромата для сохранения электронейтральности раствора, что регистрируется в виде отрицательного пика (инвертируемого в положительный в ПО «Эльфоран» для удобства обработки)~\cite{komarova2006prakticheskoe, shelekhovaOpredelenieAnionovNeorganicheskikh2023}.

\subsection{Выбор рабочего напряжения}

Напряженность электрического поля в капилляре определяет скорость движения зон компонентов и общую эффективность электрофоретической системы~\cite{komarova2006prakticheskoe}. Согласно нормативному документу ПНД Ф 14.1:2:3:4.282-18, в качестве базового установлено рабочее напряжение отрицательной полярности $-25$~кВ~\cite{pndf14123428218}. Использование высокой разности потенциалов минимизирует время разделения (до 5 минут) и снижает вклад продольной диффузии в размывание ионных зон, увеличивая пиковую емкость системы~\cite{pobozyApplicationCapillaryElectrophoresis2021b}. Однако протекание тока лимитируется тепловыделением (джоулевым теплом, $P = U \cdot I$), избыток которого вызывает флуктуации вязкости буфера, уширение пиков и нестабильность тока~\cite{komarova2006prakticheskoe}. Жидкостная система принудительного охлаждения прибора «Капель-105М» эффективно отводит тепло, обеспечивая стабильность базовой линии при $-25$~кВ~\cite{pndf14123428218}. Снижение напряжения до $-17$~кВ рекомендуется применять исключительно при анализе высокоминерализованных проб с целью предотвращения токовых перегрузок~\cite{pndf14123428218}.

\subsection{Обоснование состава ведущего электролита и расчет pH}

Состав рабочего буфера подобран с целью обеспечения селективного разделения анионов и минимизации шумов детектора~\cite{shelekhovaOpredelenieAnionovNeorganicheskikh2023}:
\begin{itemize}
	\item Оксид хрома (VI) ($\text{CrO}_3$) в концентрации 10~ммоль/дм$^3$ служит источником поглощающих хромат-ионов~\cite{pndf14123428218}. Электрофоретическая подвижность хромата близка к подвижности наиболее быстрых аналитов (хлорида и сульфата), что предотвращает электродисперсию зон и сохраняет симметрию пиков~\cite{komarova2006prakticheskoe}.
	\item Диэтаноламин (ДЭА) в концентрации 30~ммоль/дм$^3$ выступает в роли буферного компонента, создающего слабощелочную среду (рН 8,7–9,2)~\cite{pndf14123428218}. Щелочные значения необходимы для полной диссоциации фтористоводородной и фосфорной кислот до легко мигрирующих анионов $\text{F}^-$ и $\text{HPO}_4^{2-}$~\cite{komarova2006prakticheskoe, shelekhovaOpredelenieAnionovNeorganicheskikh2023}.
	\item Гидроксид цетилтриметиламмония (ЦТА-ОН) в концентрации 2~ммоль/дм$^3$ выполняет роль модификатора электроосмотического потока (ЭОП)~\cite{pndf14123428218}. Катионные молекулы ПАВ адсорбируются на отрицательно заряженных силанольных группах кварцевого стекла, формируя положительно заряженный бислой~\cite{pobozyApplicationCapillaryElectrophoresis2021b}. Это приводит к обращению направления ЭОП в сторону анода (выхода из капилляра), благодаря чему ускоряется миграция аналитов к детектору~\cite{komarova2006prakticheskoe, pobozyApplicationCapillaryElectrophoresis2021b}.
\end{itemize}

Для теоретического подтверждения стабильности pH приготовленного буфера составлена система балансовых соотношений. Диссоциация хромовой кислоты в водной среде протекает полностью по первой ступени и частично по второй ($pK_{a2} = 6,50$):
\begin{equation}
	\text{HCrO}_4^- \rightleftharpoons \text{H}^+ + \text{CrO}_4^{2-}
\end{equation}

Протонирование слабого органического основания ДЭА ($B$) описывается константой кислотности сопряженной кислоты $BH^+$, равной $pK_a \approx 8,88$ при $25~^\circ\text{C}$~\cite{komarova2006prakticheskoe}. Сильное основание ЦТА-ОН диссоциирует полностью. Уравнение электронейтральности для данной системы записывается следующим образом:
\begin{equation}
	[\text{H}^+] + [\text{CTA}^+] + [\text{BH}^+] = [\text{OH}^-] + [\text{HCrO}_4^-] + 2[\text{CrO}_4^{2-}]
\end{equation}

В щелочной области (pH 8,5–9,2) концентрации свободных ионов водорода ($[\text{H}^+] < 10^{-8}$~М) и гидроксил-ионов ($[\text{OH}^-] < 10^{-5}$~М) пренебрежимо малы по сравнению с концентрациями электролитов ($10^{-3}$–$10^{-2}$~М). Поскольку pH среды существенно выше константы $pK_{a2}$ хромовой кислоты (6,50), доля гидрохромат-ионов стремится к нулю, а концентрация двухзарядных хромат-ионов практически равна общему содержанию хрома ($[\text{CrO}_4^{2-}] \approx C_{\text{Cr}} = 0,010$~М). Баланс зарядов упрощается до линейного вида:
\begin{equation}
	[\text{CTA}^+] + [\text{BH}^+] \approx 2[\text{CrO}_4^{2-}]
\end{equation}
\begin{equation}
	0,002 + [\text{BH}^+] = 2 \cdot 0,010 = 0,020 \implies [\text{BH}^+] \approx 0,018~\text{М}
\end{equation}

Тогда концентрация неактивной формы буферного амина составляет:
\begin{equation}
	[B] = C_{\text{DEA}} - [\text{BH}^+] = 0,030 - 0,018 = 0,012~\text{М}
\end{equation}

Равновесная концентрация ионов водорода выражается через константу диссоциации протонированной формы ДЭА:
\begin{equation}
	K_a = \frac{[B][\text{H}^+]}{[\text{BH}^+]} \implies [\text{H}^+] = K_a \cdot \frac{[\text{BH}^+]}{[B]} = K_a \cdot \frac{0,018}{0,012} = 1,5 \cdot K_a
\end{equation}
\begin{equation}
	\text{pH} = pK_a - \log_{10}(1,5) = pK_a - 0,176
\end{equation}

Для стандартной температуры раствора ($25~^\circ\text{C}$, $pK_a = 8,88$) расчетное значение pH составляет $8,70$~\cite{komarova2006prakticheskoe}. С понижением рабочей температуры до $20~^\circ\text{C}$ константа $pK_a$ ДЭА возрастает до $9,00$ вследствие экзотермичности процесса протонирования аминов, что смещает теоретический pH буфера в область $8,83$~\cite{komarova2006prakticheskoe}. Полученное значение укладывается в оптимальный диапазон функционирования электролита~\cite{pndf14123428218}.

\subsection{Выбор температуры термостатирования}

Температура оказывает сильное влияние на вязкость фонового раствора и подвижность ионов~\cite{komarova2006prakticheskoe}. Термостатирование капилляра при температуре $20~^\circ\text{C}$ (с точностью регулирования до $0,1~^\circ\text{C}$) позволяет увеличить вязкость среды и снизить силу рабочего тока по сравнению с комнатными условиями, минимизируя тепловую дисперсию зон~\cite{pndf14123428218, pobozyApplicationCapillaryElectrophoresis2021b}. Это гарантирует высокую стабильность базовой линии и воспроизводимость времен миграции аналитов (относительное стандартное отклонение RSD времени выхода не превышает $1,2\%$)~\cite{shelekhovaOpredelenieAnionovNeorganicheskikh2023}.

\section{ИДЕНТИФИКАЦИЯ ПИКОВ И КАЧЕСТВЕННЫЙ АНАЛИЗ}

Идентификация индивидуальных анионов на электрофореграммах проводится по временам их миграции ($t_m$, мин)~\cite{pndf14123428218}.

\subsection{Порядок выхода компонентов}

Стандартная последовательность регистрации пиков имеет вид: хлорид $\rightarrow$ нитрит $\rightarrow$ сульфат $\rightarrow$ нитрат $\rightarrow$ фторид $\rightarrow$ фосфат~\cite{shelekhovaOpredelenieAnionovNeorganicheskikh2023}. Очередность выхода определяется величиной их эффективной электрофоретической подвижности, зависящей от соотношения ионного заряда и радиуса гидратированного иона~\cite{pobozyApplicationCapillaryElectrophoresis2021b}. Сульфат-ион ($\text{SO}_4^{2-}$), несмотря на двойной отрицательный заряд, выходит позже однозарядного нитрита ($\text{NO}_2^-$). Это обусловлено его высокой плотностью заряда, вызывающей притяжение молекул воды и формирование объемной гидратной оболочки, что увеличивает эффективный гидродинамический радиус и снижает скорость электромиграции~\cite{pobozyApplicationCapillaryElectrophoresis2021b, shelekhovaOpredelenieAnionovNeorganicheskikh2023}. Фосфаты при рН ведущего электролита находятся в форме гидрофосфата $\text{HPO}_4^{2-}$, обладающего меньшей подвижностью и детектируемого последним~\cite{komarova2006prakticheskoe}.

\subsection{Метод стандартных добавок}

Для надежной качественной идентификации анионов в пробах со сложным солевым фоном применяли метод стандартных добавок~\cite{pndf14123428218}. В анализируемую воду вносили микроколичества смесей ГСО целевых анионов~\cite{shelekhovaOpredelenieAnionovNeorganicheskikh2023}. Пик идентифицировали по симметричному приращению его высоты и площади при сохранении формы зоны и отсутствии расщепления сигнала~\cite{pndf14123428218}. Низкие значения RSD времени миграции (до $1,2\%$) подтверждают воспроизводимость условий электрофореза.

\section{МЕТРОЛОГИЧЕСКИЕ ХАРАКТЕРИСТИКИ И ГРАДУИРОВКА}

Количественная оценка концентраций анионов основана на градуировочных зависимостях площадей пиков ($S$, отн. ед.) от концентрации ($C$, мг/дм$^3$)~\cite{pndf14123428218}. Градуировочные характеристики рассчитывались в ПО «Эльфоран» по методу наименьших квадратов~\cite{shelekhovaOpredelenieAnionovNeorganicheskikh2023}.

\subsection{Построение градуировочных графиков}

Уравнения регрессии аппроксимированы в виде линейных функций вида $C = k_1 \cdot S$, проходящих через начало координат~\cite{shelekhovaOpredelenieAnionovNeorganicheskikh2023}. Метрологические характеристики градуировочных моделей представлены в таблице \ref{tab:metro_params}.

\begin{table}[H]
\centering
\caption{Метрологические и градуировочные характеристики определения анионов}
\label{tab:metro_params}
\begin{tabular}{|l|c|c|c|c|c|}
\hline
Анион & Диапазон, мг/дм$^3$ & $t_m$, мин & Уравнение градуировки & $R^2$ & LOD / LOQ, мг/дм$^3$ \\ \hline
Хлорид   & 0,5–200,0 & 3,35 & $C = 0,1248 \cdot S$ & 0,9998 & 0,15 / 0,50 \\ \hline
Нитрит   & 0,2–50,0  & 3,50 & $C = 0,1600 \cdot S$ & 0,9998 & 0,06 / 0,20 \\ \hline
Сульфат  & 0,5–200,0 & 3,59 & $C = 0,1531 \cdot S$ & 0,9998 & 0,15 / 0,50 \\ \hline
Нитрат   & 0,2–50,0  & 3,77 & $C = 0,2004 \cdot S$ & 0,9998 & 0,06 / 0,20 \\ \hline
Фторид   & 0,1–10,0  & 4,10 & $C = 0,0445 \cdot S$ & 0,9994 & 0,03 / 0,10 \\ \hline
Фосфат   & 0,25–25,0 & 4,44 & $C = 0,1009 \cdot S$ & 0,9997 & 0,08 / 0,25 \\ \hline
\end{tabular}
\end{table}

Высокие коэффициенты детерминации ($R^2 \ge 0,999$) подтверждают корректность линейного приближения. В области концентраций выше 200 мг/дм$^3$ наблюдается отклонение от линейности вследствие электродисперсии зон из-за различия проводимости зоны аналита и фонового буфера, что приводит к выраженной асимметрии пиков~\cite{komarova2006prakticheskoe}. Для анализа таких проб требуется предварительное разбавление деионизованной водой~\cite{pndf14123428218}.

\subsection{Чувствительность и воспроизводимость}

Пределы обнаружения (LOD) оценивались по критерию тройного превышения аналитического сигнала над уровнем шума ($S/N = 3$), пределы количественного определения (LOQ) — по соотношению $S/N = 10$~\cite{filiposkaSpasicValidationIon2025}. Оценка сходимости результатов параллельных вводов показала, что относительное стандартное отклонение RSD площадей пиков не превышает $5,0\%$, что соответствует требованиям ГОСТ Р ИСО 5725-2002~\cite{pndf14123428218}.

\section{АНАЛИЗ РЕАЛЬНОГО ОБЪЕКТА И ОБСУЖДЕНИЕ РЕЗУЛЬТАТОВ}

Для практической апробации электрофоретической схемы выполнен анализ серии параллельных проб поверхностной воды из реки Охта (Ленинградская область)~\cite{shelekhovaOpredelenieAnionovNeorganicheskikh2023}.

\subsection{Результаты количественного химического анализа}

Статистическая обработка результатов измерений шести параллельных образцов («галя а1» – «галя а9» из папки данных) проведена при доверительной вероятности $P = 0,95$. Результаты количественного анализа и их сопоставление с нормативами предельно допустимых концентраций (ПДК) приведены в таблице \ref{tab:ohta_results}.

\begin{table}[H]
\centering
\caption{Результаты определения анионов в воде р. Охта и их сопоставление с ПДК}
\label{tab:ohta_results}
\begin{tabular}{|l|c|c|c|c|}
\hline
Анион & Найдено ($\bar{X} \pm \Delta X$), мг/дм$^3$ & RSD, \% & ПДК СанПиН~\cite{SanPiN368521} & ПДК Росрыбхоз~\cite{OrderRosrybolovstvo296} \\ \hline
Хлорид   & $25,57 \pm 15,32$ ($19,80 \pm 4,83$)* & 57,1 (20,9)* & 350 & 300 \\ \hline
Нитрит   & $0,031 \pm 0,010$**                   & 32,0         & 3,0 & 0,08 \\ \hline
Сульфат  & $29,52 \pm 3,10$                      & 10,0         & 500 & 100 \\ \hline
Нитрат   & $1,03 \pm 0,25$                       & 23,4         & 45 & 40 \\ \hline
Фторид   & $< 0,10$ (не обнаружено)              & —            & 1,5 & 0,05 \\ \hline
Фосфат   & $0,418 \pm 0,041$                     & 9,3          & 3,5 & 0,15 \\ \hline
\end{tabular}
\newline
\raggedright
\small
* В скобках указаны значения при исключении пробы «а9» как локальной аномалии/выброса.\\
** Значение для проб, где концентрация превышала предел обнаружения.
\end{table}

\subsection{Обсуждение экологического состояния водного объекта}

Анализ экспериментальных данных позволяет сделать экологические выводы о состоянии контролируемого участка реки Охта:
\begin{itemize}
	\item Содержание хлоридов ($25,57$~мг/дм$^3$) и сульфатов ($29,52$~мг/дм$^3$) существенно ниже установленных лимитов как для хозяйственно-бытового, так и для рыбохозяйственного водопользования~\cite{SanPiN368521, OrderRosrybolovstvo296}. Это свидетельствует об отсутствии галогенного загрязнения и стабильном солевом балансе водоема~\cite{kaushalFreshwaterSalinizationSyndrome2021, wuReviewChlorideIon2021}.
	\item Концентрации нитрит- и нитрат-ионов не превышают ПДК~\cite{SanPiN368521, OrderRosrybolovstvo296}. Это указывает на низкий уровень азотного загрязнения и эффективное самоочищение водного объекта от органических веществ~\cite{saalidongExaminingDynamicsRelationship2022}.
	\item Содержание фторидов находится ниже предела количественного определения ($<0,10$~мг/дм$^3$), что удовлетворяет санитарным правилам~\cite{SanPiN368521}. Однако LOQ метода ($0,10$~мг/дм$^3$) выше рыбохозяйственного норматива ПДК ($0,05$~мг/дм$^3$)~\cite{OrderRosrybolovstvo296}. Это отражает системное ограничение капиллярного электрофореза при контроле следовых количеств фторидов в рыбохозяйственных водоемах без предварительного концентрирования пробы~\cite{komarova2006prakticheskoe}.
	\item Обнаружено превышение рыбохозяйственного лимита по фосфатам в $2,8$ раза ($0,418$~мг/дм$^3$ при ПДК $0,15$~мг/дм$^3$)~\cite{OrderRosrybolovstvo296}. Превышение содержания фосфора является фактором эвтрофикации водоема, обусловленным стоком минеральных удобрений с сельхозугодий или сбросом неочищенных хозяйственно-бытовых сточных вод с моющими средствами~\cite{marinRelativeEffectsEutrophication2025, qinEutrophicationControlLarge2023}. Это стимулирует неконтролируемое развитие сине-зеленых водорослей и может вызвать локальный дефицит кислорода, опасный для ихтиофауны~\cite{marinRelativeEffectsEutrophication2025}.
\end{itemize}

Высокая сходимость параллельных результатов ($RSD \le 10,0\%$ для большинства ионов) подтверждает надежность предложенной схемы КЗЭ на приборах «Капель-105М» для оперативного экологического мониторинга природных вод.
